\subsection{Enriched categories}

\bdefn[title={$\A$-enriched category}]

    Let $\A$ be a monoidal category with unit $\1$.
    An {\emphcolor $\A$-enriched category} $\C$ consists of the following data:
    \benum
        \item A collection $\Ob(\C)$, whose elements we refer to as {\emphcolor objects} of $\C$.
        We abuse notation by writing $X\in\C$ instead of $X\in\Ob(\C)$.
        \item For every $X,Y\in\C$ an object of $\A$, $\ehom_\C(X,Y)$.
        \item For every $X,Y,Z\in\C$ an $\A$-morphism
        $$ c_{Z,Y,X}\colon\ehom_\C(Y,Z)\otimes\ehom_\C(X,Y)\to\ehom_\C(X,Z) $$
        which we call {\emphcolor composition}.
        \item For every $X\in\C$ a morphism $e_X\colon\1\to\ehom_\C(X,X)$ called the {\emphcolor identity} of $X$.
    \eenum
    The data must satisfy the following conditions:
    \benum
        \item Associativity: for $W,X,Y,Z\in\C$:
        \diagram{
            \ehom(Y,Z)\otimes(\ehom(X,Y)\otimes\ehom(W,X))&\ehom(Y,Z)\otimes\ehom(W,Y)\cr
            &\ehom(W,Z)\cr
            (\ehom(Y,Z)\otimes\ehom(X,Y))\otimes\ehom(W,X)&\ehom(X,Z)\otimes\ehom(W,X)
        }{
            \diagarrow{from={1,1}, to={1,2}, text={\id\otimes c_{Y,X,W}}, y distance=.25cm}
            \diagarrow{from={3,1}, to={3,2}, text={c_{Z,Y,X}\otimes\id}, y distance=-.25cm}
            \diagarrow{from={1,1}, to={3,1}, text=\alpha, x distance=-.25cm}
            \diagarrow{from={1,2}, to={2,2}, text={c_{Z,Y,W}}, x distance=.25cm}
            \diagarrow{from={3,2}, to={2,2}, text={c_{Z,X,W}}, x distance=.25cm}
        }
        Here $\alpha$ is the associator of $\A$.

        \item Unit: for every $X,Y\in\C$ the two diagrams below must commute:

        \diagram{
            \1\otimes\ehom(X,Y)&\ehom(Y,Y)\otimes\ehom(X,Y)\cr
            &\ehom(X,Y)
        }{
            \diagarrow{from={1,1}, to={1,2}, text={e_Y\otimes\id}, y distance=.25cm}
            \diagarrow{from={1,2}, to={2,2}, text={c_{Y,Y,X}}, x distance=.25cm}
            \diagarrow{from={1,1}, to={2,2}, text=\lambda, x distance=-.25cm}
        }
        \diagram{
            \ehom(X,Y)\otimes\1&\ehom(X,Y)\otimes\ehom(X,X)\cr
            &\ehom(X,Y)
        }{
            \diagarrow{from={1,1}, to={1,2}, text={\id\otimes e_X}, y distance=.25cm}
            \diagarrow{from={1,2}, to={2,2}, text={c_{Y,X,X}}, x distance=.25cm}
            \diagarrow{from={1,1}, to={2,2}, text=\rho, x distance=-.25cm}
        }

        Here $\lambda,\rho$ are the unitors of $\A$.
    \eenum

\edefn

\bexam

    $\Set$-enriched categories are precisely the locally small categories.

\eexam

\bexam

    Let $\A$ be a monoidal category and $\C$ an $\A$-enriched category.
    For every object $X\in\C$, the composition law and unit
    $$ c_{X,X,X} \colon \ehom(X,X)\otimes\ehom(X,X) \to \ehom(X,X),\qquad e_X\colon\1\to\ehom(X,X) $$
    turn $\ehom(X,X)$ into a monoid object of $\A$.
    This induces a bijection between $\Mon(\A)$ and $\A$-enriched categories with a single object.

    Thus enriched category theory can be seen as a generalization of the theory of monoid objects.

\eexam

\bremark[title={Functorality}, name={funcencat}]

    Let $\A,\A'$ be monoidal categories and $F\colon\A\to\A'$ a lax monoidal functor, with tensoror $\mu$ and unit $\epsilon$.
    Then every $\A$-enriched category $\C$ defines an $\A'$-enriched category $\C'$ as follows:
    \benum
        \item The objects of $\C'$ are the same as those of $\C$'s: $\Ob(\C')=\Ob(\C)$.
        \item For $X,Y\in\C'$ set $\ehom_{\C'}(X,Y)=F\ehom_\C(X,Y)$.
        \item For $X,Y,Z\in\C'$ define the composition law $c'_{X,Y,Z}$ in $\C'$ as the composition
        $$ \vbox{\halign{\hfil$#$\tabskip=.25cm&\hfil$#$\hfil&$#$\hfil\tabskip=0pt\cr
            \ehom_{\C'}(Y,Z)\otimes\ehom_{\C'}(X,Y) &=& F(\ehom_\C(Y,Z))\otimes F(\ehom_\C(X,Y))\cr
                &\xtos\mu& F(\ehom_\C(Y,Z)\otimes\ehom_\C(X,Y))\cr
                &\xtos{Fc_{Z,Y,X}}& F(\ehom_\C(X,Z))\cr
                &=&\ehom_{\C'}(X,Z)\cr
        }} $$
        \item For every $X\in\C'$ the identity $e'_X$ is the composition
        $$ \1_{\A'} \xtos\epsilon F(\1_\A) \xtos{F(e_X)} F(\ehom_\C(X,X)) = \ehom_{\C'}(X,X) . $$
    \eenum

\eremark

\bdefn[title={The underlying category}, name={underencat}]

    Let $\A$ be a monoidal category and $\C$ an $\A$-enriched category.
    Consider then the monoidal functor $F\colon\A\to\Set$ given by $F=\hom_\A(\1,\bul)$ (\refmath[example]{laxcorepfunc}).
    By the above remark, this gives rise to a $\Set$-enriched category, which we can identify with an ordinary category.
    We call this the {\emphcolor underlying category} of $\C$, and we abuse notation and denote it by $\C$.

    Concretely, this category has the same objects as $\C$, and its morphisms are
    $$ \hom_\C(X,Y) = \hom_\A(\1,\ehom_\C(X,Y)) . $$

\edefn

\bnote

    Let $\A$ be a monoidal category and $\C$ an ordinary category.
    An {\emphcolor $\A$-enrichment} of $\C$ is an $\A$-enriched category $\hat\C$ such that the underlying category of $\hat\C$ is isomorphic to
    $\C$.

\enote

\bexam

    Let $\Top$ denote the category of topological spaces, endowed with the monoidal structure given by the cartesian product.
    A $\Top$-enriched category is called a {\emphcolor topologically enriched category}.
    Note that the corepresentable functor $\hom_\Top(\1,\bul)$ (where $\1=\set*$ is the unit/terminal object of $\Top$) is naturally isomorphic
    to the forgetful functor $\Top\to\Set$.

    Thus is $\C$ is a topologically enriched category, its underlying category $\C_0$ can be described as follows:
    \benum
        \item Its objects are the same as the objects of $\C$.
        \item For $X,Y\in\C_0$, morphisms $X\to Y$ in $\C_0$ are the same as points in the space $\ehom_\C(X,Y)$.
        \item Given $f\colon X\to Y$ and $g\colon Y\to Z$ in $\C_0$, the composition $g\circ f$ is the image of $(g,f)$ under the continuous map
        $$ c_{Z,Y,X}\colon\ehom_\C(Y,Z)\times\ehom_\C(X,Y) \to \ehom_\C(X,Z) . $$
        \item For $X\in\C_0$, $\id_X$ is the image of $*\in\set*$ under $e_X\colon\set*\to\ehom_\C(X,X)$.
    \eenum

    Conversely, enriching an ordinary category $\C_0$ over $\Top$ is the same as endowing each hom-set $\hom_{\C_0}(X,Y)$ with a topology such that
    composition is continuous $\hom_{\C_0}(Y,Z)\times\hom_{\C_0}(X,Y)\to\hom_{\C_0}(X,Z)$.

\eexam

\bnote

    The identities of an enriched category are unique.
    That is, if $\C,\C'$ are two $\A$-enriched categories with the same objects, morphisms, and compositions, then $e_X=e'_X$ for all $X\in\C$.
    Indeed, $e_X$ is simply the unit of the $\A$-monoid $\ehom_\C(X,X)$ (with $c_{X,X,X}$ as multiplication) which we showed is unique.

\enote

\bdefn[title={Enriched functors}]

    Let $\A$ be a monoidal category and $\C,\D$ be $\A$-enriched categories.
    An $\A$-enriched functor $F\colon\C\to\D$ consists of the following data:
    \benum
        \item For every $X\in\C$ an associated object $F(X)\in\D$.
        \item For every $X,Y\in\C$ a morphism in $\A$
        $$ F_{X,Y}\colon\ehom_\C(X,Y) \to \ehom_\D(F(X),F(Y)) . $$
    \eenum
    This data is subject to the following conditions:
    \benum
        \item For every $X\in\C$, the morphism $e_{FX}\colon\1\to\ehom_\D(X,X)$ factors as
        $$ \1 \xtos{e_X} \ehom_\C(X,X) \xtos{F_{X,X}} \ehom_\D(X,X) . $$
        \item For every $X,Y,Z\in\C$ the following diagram commutes

        \diagram{
            \ehom_\C(Y,Z)\otimes\ehom_\C(X,Y)&\ehom_\C(X,Z)\cr
            \ehom_\D(F(Y),F(Z))\otimes\ehom_\D(F(X),F(Y))&\ehom_\D(F(X),F(Z))
        }{
            \diagarrow{from={1,1}, to={1,2}, text={c_{Z,Y,X}}, y distance=.25cm}
            \diagarrow{from={2,1}, to={2,2}, text={c_{FZ,FY,FX}}, y distance=-.25cm}
            \diagarrow{from={1,1}, to={2,1}, text={F_{Y,Z}\otimes F_{X,Y}}, x distance=-.25cm}
            \diagarrow{from={1,2}, to={2,2}, text={F_{X,Z}}, x distance=.25cm}
        }
    \eenum

    If $\C\xtos F\D\xtos G\E$ are $\A$-enriched functors, then we can define their composition $G\circ F\colon\C\to\E$ by
    $$ (G\circ F)(X) = G(F(X)),\ X\in\C,\qquad (G\circ F)_{X,Y} = G_{FX,FY}\circ F_{X,Y} . $$

    Thus we can form the (ordinary) category $\Cat(\A)$ of small (meaning $\Ob(\C)$ is small) $\A$-enriched categories, whose morphisms
    are $\A$-enriched functors.

\edefn

\bexam

    A $\Set$-enriched functor is easily seen to be an ordinary functor.
    Thus there is an isomorphism of categories $\Cat\cong\Cat(\Set)$.

\eexam

\bexam

    If $F\colon\A\to\A'$ is a lax monoidal functor between monoidal categories, then we showed that it maps $\A$-enriched categories $\C$
    to $\A'$-enriched categories $\C'$.
    Similarly it maps an $\A$-enriched functor $f\colon\C\to\D$ to the $\A'$-enriched functor $f'\colon\C'\to\D'$ which acts on objects
    as $f$ does, and $f'_{X,Y}=F(f_{X,Y})$.

    This defines a functor of ordinary categories $\Cat(\A)\to\Cat(\A')$.
    In particular for the functor $F=\ehom_\A(\1,\bul)$, we obtain a forgetful functor
    $$ \Cat(\A) \to \Cat(\Set) \cong \Cat $$
    which assigns to each (small) $\A$-enriched category its underlying ordinary category.

\eexam

We can also define enriched natural transformations.
Recall that an ordinary natural transformation $\alpha\colon F\To G\colon\C\to\D$ is given by a $\C$-indexed family of $\D$-morphisms
$\alpha_X\colon FX\to GX$ such that the usual diagram commutes.
We can write this naturality condition as requiring that the following square commutes for all $X,Y\in\C$:

\commsq{
    \hom_\C(X,Y)&\hom_\D(FX,FY)\cr
    \hom_\D(GX,GY)&\hom_\D(FX,GY)
}{F}{\circ\alpha_X}{G}{\alpha_Y\circ}

In enriched category theory in order to talk about ``picking out'' a morphism $\alpha_X$, we talk instead about a morphism
$\alpha_X\colon\1\to\ehom_\D(FX,FY)$.
This motivates the following.

\bdefn[title={Enriched natural transformations}]

    Let $F,G\colon\C\to\D$ be $\A$-enriched functors.
    An {\emphcolor $\A$-enriched natural transformation} $\alpha\colon F\To G$ is an $\Ob(\C)$ indexed family of $\A$-morphisms
    $\set{\alpha_X\colon\1\to\ehom_\D(FX,GX)}_{X\in\C}$ such that the following diagram commutes (in $\A$) for all $X,Y\in\C$:

    \diagram{
        \1\otimes\ehom_\C(X,Y)&\ehom_\D(FY,GY)\otimes\ehom_\D(FX,FY)\cr
        \ehom_\C(X,Y)&\ehom_\D(FX,GY)\cr
        \ehom_\C(X,Y)\otimes\1&\ehom_\D(GX,GY)\otimes\ehom_\D(FX,GX)
    }{
        \diagarrow{from={2,1}, to={1,1}, text={\lambda^{-1}}, x distance=-.25cm}
        \diagarrow{from={2,1}, to={3,1}, text={\rho^{-1}}, x distance=-.25cm}
        \diagarrow{from={1,1}, to={1,2}, text={\alpha_Y\otimes F}, y distance=.25cm}
        \diagarrow{from={3,1}, to={3,2}, text={G\otimes\alpha_X}, y distance=-.25cm}
        \diagarrow{from={1,2}, to={2,2}, text=c, x distance=.25cm}
        \diagarrow{from={3,2}, to={2,2}, text=c, x distance=.25cm}
    }

\edefn

As with ordinary natural transformations, we can compose enriched natural transformations both vertically and horizontally.

\bdefn

    Suppose $\C,\D$ are $\A$-enriched categories and we have enriched functors and natural transformations as in the following diagram:

    \diagram{
        &\cr
        \C&&\D\cr
        &
    }{
        \diagarrow{from={2,1}, to={2,3}, curve=3cm, text=F, y distance=1.5cm, y off=.1cm}
        \diagarrow{from={2,1}, to={2,3}, curve=-3cm, text=H, y distance=-1.5cm, y off=-.1cm}
        \diagarrow{from={2,1}, to={2,3}, text=G, y distance=.25cm, x distance=.25cm}
        \diagarrow{from={1,2}, to={2,2}, text=\alpha, ds, y off=-.25cm, x distance=.25cm}
        \diagarrow{from={3,2}, to={2,2}, text=\beta, ds, y off=.25cm, x distance=.25cm, left cap=<, right cap=-}
    }

    Then we define the {\emphcolor vertical composition} $\beta\circ\alpha\colon F\To H$ to be
    $$ (\beta\circ\alpha)_X = \1 \xtos{v^{-1}} \1\otimes\1 \xtos{\beta_X\otimes\alpha_X} \ehom_\D(GX,HX)\otimes\ehom_\D(FX,GX) \xtos c
    \ehom_\D(FX,HX) . $$
    And if we have a diagram of the form

    {\def\diagrowheight{0pt}\def\diagrowbuf{0pt}\diagram{
        &&\cr
        \C&\C'&&\D'&\D\cr
        &&
    }{
        \diagarrow{from={2,1}, to={2,2}, text=F, y distance=.25cm}
        \diagarrow{from={2,4}, to={2,5}, text=G, y distance=.25cm}
        \diagarrow{from={2,2}, to={2,4}, curve=1cm, y off=.1cm, text={H}, y distance=.75cm}
        \diagarrow{from={2,2}, to={2,4}, curve=-1cm, y off=-.1cm, text={K}, y distance=-.75cm}
        \diagarrow{from={1,3}, to={3,3}, ds, text=\alpha, x distance=.25cm}
    }}

    We define the {\emphcolor whiskering} of $\alpha$ to be
    $$ \displaylines{
        (G\alpha)_X = \1 \xtos{\alpha_X} \ehom_{\D'}(HX,KX) \xtos G \ehom_{\D}(GHX,GKX) \colon G\circ H\To G\circ K\cr
        (\alpha F)_X = \alpha_{FX} \colon H\circ F\To K\circ F
    } $$

    Finally if we a diagram of the form

    {\def\diagrowheight{0pt}\def\diagrowbuf{0pt}\diagram{
        &&&\cr
        \C&&\D&&\E\cr
        &&&
    }{
        \diagarrow{from={2,1}, to={2,3}, curve=1cm, y off=.1cm, text={F}, y distance=.75cm}
        \diagarrow{from={2,1}, to={2,3}, curve=-1cm, y off=-.1cm, text={G}, y distance=-.75cm}
        \diagarrow{from={2,3}, to={2,5}, curve=1cm, y off=.1cm, text={H}, y distance=.75cm}
        \diagarrow{from={2,3}, to={2,5}, curve=-1cm, y off=-.1cm, text={K}, y distance=-.75cm}
        \diagarrow{from={1,2}, to={3,2}, ds, text=\alpha, x distance=.25cm}
        \diagarrow{from={1,4}, to={3,4}, ds, text=\beta, x distance=.25cm}
    }}

    We define the {\emphcolor horizontal composition} $\beta\star\alpha\colon H\circ F\To K\circ G$ to be
    $$ \beta\star\alpha = (\beta G)\circ(H\alpha) = (K\alpha)\circ(\beta F) . \eqcnt $$

\edefn

It is straightforward to show that vertical composition, whiskering, and horizontal composition are all natural transformations.
It is furthermore straightforward to show the identity \gotoleqn.
Furthermore it is straightforward to show that $\Cat(\A)$ forms a strict $2$-category (i.e.\ a category enriched over $\Cat$).

\bremark

    Recall the construction of an underlying category for an $\A$-enriched category $\C$, denoted $\C_0$.
    As it turns out this construction is functorial, $2$-functorial even.

    Let $\I$ be the {\emphcolor unit $\A$-category} with a single object $*$ and $\I(*,*)=\1$; composition is given by the unit $v$.
    Then the underlying category functor is just the representable $2$-functor
    $$ \bul_0 \colon \Cat(\A) \to \Cat,\qquad \bul_0 = \hom_{\Cat(\A)}(\I,\bul) . $$
    Indeed, $\A$-functors $\I\to\C$ correspond to a choice of object $A\in\C$.
    A natural transformation $A\To B\colon\I\to\C$ corresponds to a morphism $\1\to\ehom_\C(A,B)$.
    Thus the category $\hom_{\Cat(\A)}(\I,\C)$ has as objects the objects of $\C$ and as morphisms the morphisms $\1\to\ehom_\C(A,B)$
    for $A,B\in\C$.
    This is precisely $\C_0$.

    Note that we can write this as
    $$ \hom_{\C_0}(A,B) = \A_\1(\ehom_\C(A,B)), $$
    where $\A_\1\colon\A\to\Set$ is the representable functor $\A_\1=\hom_\A(\1,\bul)$.
    It then follows that the relation of $\C_0$ and $\C$ depends heavily on how faithful the functor $\A_\1$ is.

    For example if $\A=\Cat$ then $\Cat_\1$ is not faithful: $\C$ is a $2$-category while $\C_0$ is the category obtained by forgetting the
    $2$-cells.
    When $\A=\Top$ then $\C_0$ only loses the topology on the hom-sets.
    And when $\A=\Ab$ or $\Mod_R$ then $\A_\1$ is faithful and conservative (reflects isomorphisms) so that $\C_0$ is relatively similar to $\C$.

    An enriched functor $F\colon\C\to\D$ then corresponds to the ordinary functor $F_0\colon\C_0\to\D_0$ which maps $A\colon\I\to\C$ to
    $F\circ A$ and $f\colon A\to B$ to the whiskering $Ff$, which is the composite
    $$ \1 \xtos f\ehom_\C(A,B) \xtos F \ehom_\D(FA,FB) . $$
    As we identify $A\colon\I\to\C$ with $A\in\C$, and $Ff\colon\1\to\ehom_\D(FA,FB)$ with morphisms $FA\to FB$ this tells us that
    $$ F_0(A) = FA,\qquad F_0(f) = Ff . $$
    In other words,
    $$ (F_0)_{A,B} \colon \hom_{\C_0}(A,B) \to \hom_{\D_0}(FA,FB) \hbox{ is } \A_\1 F_{A,B} \colon \A_1\ehom_\C(A,B) \to \A_\1\ehom_\D(FA,FB) . $$

    For $\A$-natural transformations $\alpha\colon F\To G$, these become ordinary natural transformations $\alpha_0\colon F_0\To G_0$
    whose components $(\alpha_0)_A\in\hom_{\D_0}(FA,GA)$ is precisely $\alpha_A\colon\1\to\ehom_\D(FA,GA)$.

\eremark

